\scanpage{39}
% slanted
\margin{Qui sont $\mathcal A,\mathcal B$ ?}
\deleted{\unsure{$\ell$}} 3.1.7. Pour $\underline F$ \deleted{complexe des}\inserted{un} $f^*\mathcal B$-Module\deleted{s} à droite sur $X$ et $\underline G$ \deleted{le complexe \unsure{de}} un $\mathcal A,\mathcal B$-Bimodule à gauche\inserted{sur $S$}, on dispose de
\[
 f_!^\bullet\underline F\otimes_{\mathcal B}\underline G
 \longrightarrow f_!^\bullet(\underline F\otimes_{\mathcal B} f^*\underline G).
 \tag*{(3.1.7.1)}
\]
Pour $\underline H$ faisceau de $\mathcal A$-modules sur $S$, cette flèche se transpose en
\[
 \underline{\mathrm{Hom}}_{\mathcal A}(f^*\underline G,f_\bullet^!\underline H)
 \longrightarrow f_\bullet^!\underline{\mathrm{Hom}}_{\mathcal A}(\underline G,\underline H)
 \tag*{(3.1.7.2)}
\]
Cette flèche se dérive en une flèche \deleted{d'induction}
\[
 R\underline{\mathrm{Hom}}_{\mathcal A}(f^*\underline{\deleted{\unsure{G}}L},Rf^!\underline{\deleted{\unsure{H}}M})
 \xrightarrow{\sim} Rf^!R\underline{\mathrm{Hom}}_{\mathcal A}(\underline{\deleted{\unsure{G}}L},\underline{\deleted{\unsure{H}}M}),
 \tag*{(3.1.7.3)}
\]
\margin{\tikz{\draw (.07,0) -- (0,.07) -- (.08,.43); \draw (0,.07) -- (-.02,0);}}
\inserted{($L\in\mathrm{ob}$ ?, $M\in\mathrm{ob}$ ?), appelée \uline{homomorphisme} \uline{d'induction}, dont on va vérifier que c'est un isomorphisme. Il} \deleted{qui} rend commutatif le diagramme suivant, pour $K\in D^-(X,f^*\mathcal B)$, $L\in D^-(S,\mathcal A,\mathcal B)$
\margin{\unsure{réf. not.|réf.} ?}
\deleted{$\mathrm{Hom}(Rf_!K,R\underline{\mathrm{Hom}}_{\mathcal A}(f^*L,Rf^!M))$} et $M\in D^+(S,\mathcal A)$.
\diagram{\resizebox{\linewidth}{!}{\begin{tikzcd}[ampersand replacement=\&,column sep=1.3em,row sep=2.5em]
 \mathrm{Hom}(K,R\underline{\mathrm{Hom}}_{\mathcal A}(f^*L,Rf^!M))
 \arrow[r,"\text{\textcircled{\scriptsize 1}}"] \arrow[d,leftrightarrow,"\sim" description]
 \& \mathrm{Hom}(K,Rf^!R\underline{\mathrm{Hom}}_{\mathcal A}(L,M))
 \arrow[r,leftrightarrow,"\sim","\mathrm{adj.}"']
 \& \mathrm{Hom}(Rf_!K,R\underline{\mathrm{Hom}}_{\mathcal A}(L,M))
 \arrow[d,leftrightarrow,"\sim" description]\\
 \mathrm{Hom}(K\overset{\mathbb L}{\otimes}_{\mathcal B}f^*L,Rf^!M)
 \arrow[r,leftrightarrow,"\sim","\mathrm{adj}"']
 \& \mathrm{Hom}(Rf_!(K\overset{\unsure{L|4}}{\otimes}f^*L),M)
 \arrow[r,"\text{\textcircled{\scriptsize 2}}"]
 \& \mathrm{Hom}(Rf_!K\overset{\mathbb L}{\otimes}L,M)
\end{tikzcd}}}
% slanted
\margin{terminologie}
Dans ce diagramme, \textcircled{\scriptsize 2} (déduit de la flèche de \deleted{changement de base}\inserted{projection}) est un isomorphisme ($\overline{\underline{\mathrm{XVII}}}$ 5.2). La flèche \textcircled{\scriptsize 1} est donc un isomorphisme quel que soit $K$. On vérifie facilement que ceci implique que (3.1.7.3) est un isomorphisme. \inserted{On a obtenu}

\margin{Préciser données $f$ comp. $X\to S$, $\mathcal A$ $\mathcal B$ \unsure{(?)}}
\uline{Proposition} 3.1.8. \uline{La flèche} 3.1.7.3 \uline{est un isomorphisme} : \uline{on a la ``\uwave{formule d'induction}''} :
\[
 R\underline{\mathrm{Hom}}_{\mathcal A}(f^*L,Rf^!M)
 \xrightarrow{\sim} Rf^!R\mathrm{Hom}_{\mathcal A}(L,M).
\]
\inserted{pour $L\in\mathrm{ob}\qquad,\ M\in\mathrm{ob}\qquad$.}
% archive sheet number: 38
